\section{Contextualização}
\label{sec:1.1}

A expansão do ensino superior nas últimas décadas intensificou a complexidade da gestão acadêmica nas Instituições de Ensino Superior (IES). O crescimento do número de cursos, turmas, professores e disciplinas ofertadas por semestre amplia significativamente as combinações possíveis de alocação, tornando a organização de grades horárias uma tarefa marcada por forte crescimento combinatório. Nesse contexto, o \textit{University Course Timetabling Problem} (UCTP) emerge como um dos desafios logísticos mais recorrentes para coordenadores acadêmicos: distribuir eventos de ensino, como aulas, laboratórios e atividades práticas, em intervalos de tempo disponíveis, respeitando simultaneamente as disponibilidades docentes, os vínculos curriculares e as restrições acadêmicas de cada curso \cite{ceschia2023}.

O problema se estrutura em torno de dois tipos de restrições. As restrições rígidas (\textit{hard constraints}) são obrigatórias, de modo que sua violação compromete a validade da grade produzida. Exemplos típicos são a proibição de alocar um mesmo professor em dois horários simultâneos e a impossibilidade de atribuir uma mesma turma a duas atividades no mesmo período. As restrições flexíveis (\textit{soft constraints}), por sua vez, não invalidam necessariamente a solução, mas afetam sua qualidade. A preferência de um docente por determinado turno, a redução de janelas ociosas e a distribuição mais equilibrada da carga horária semanal são exemplos de critérios que, quando ignorados, podem degradar a experiência de professores e estudantes \cite{weare1995}. A satisfação simultânea dessas categorias, em instâncias de maior porte, ultrapassa rapidamente a viabilidade do planejamento manual.

À medida que o número de professores, turmas, disciplinas e horários cresce, o espaço de busca sofre uma explosão combinatória que inviabiliza tanto a resolução manual quanto a enumeração exaustiva por métodos computacionais. Métodos exatos, embora relevantes em instâncias menores ou altamente controladas, tendem a apresentar limitações de escalabilidade diante dos volumes característicos de uma universidade de médio porte \cite{bashab2023}. Essa intratabilidade justifica a adoção de abordagens heurísticas e meta-heurísticas, que sacrificam a garantia de otimalidade em troca de soluções de boa qualidade obtidas em tempo computacional aceitável.

Diante desse cenário, o presente trabalho propõe o desenvolvimento do IGRIS (\textit{Intelligent Grade Integration System}), um sistema de suporte à decisão voltado à geração automatizada de grades horárias acadêmicas por meio de Algoritmos Genéticos (AGs). O IGRIS adota uma arquitetura de monolito modular com adaptadores de microsserviços, integrando três implementações independentes do algoritmo genético desenvolvidas em paradigmas de programação distintos: orientado a objetos, com o Sunny, desenvolvido em Java com Spring Boot; funcional, com o Morpheus, desenvolvido em Python com Flask; e concorrente, com o Rocket, desenvolvido em Go com \textit{goroutines}. Essa estrutura permite não apenas automatizar a geração de grades, mas também comparar quantitativamente o desempenho de diferentes abordagens computacionais aplicadas ao mesmo problema de otimização.
